% ============================================================================
%  CAPÍTULO II: ALCANCE Y LÍMITES
% ============================================================================

\section{Alcance y Justificación}

\subsection{Justificación}

La implementación de mecanismos de defensa en profundidad representa una necesidad operativa imperiosa para el sostenimiento de cualquier infraestructura corporativa. En entornos donde un único servidor centralizado (\texttt{dc1}) asume múltiples funciones neurálgicas (Controlador de Dominio Active Directory, DNS autoritativo, RADIUS para acceso inalámbrico 802.1X, almacenamiento de archivos departamentales y pasarela web), cualquier intrusión no detectada o exploit no contenido tiene el potencial de comprometer la totalidad de la organización.

El presente estudio se justifica en la necesidad de dotar a los administradores de sistemas de herramientas de monitoreo capaces de operar en tiempo de ejecución, ofreciendo tanto la visibilidad forense y telemetría de eventos requerida para el análisis de incidentes (IDS) como la capacidad de mitigación proactiva y automática de ataques directos (IPS).

\subsection{Alcance Técnico}

El alcance del proyecto comprende las siguientes dimensiones de trabajo:

\begin{itemize}
  \item \textbf{Dimensión Teórica y Comparativa:} Análisis conceptual exhaustivo de las diferencias operativas y riesgos entre IDS e IPS, acompañado de una evaluación técnica pormenorizada de diez criterios de ingeniería entre las plataformas \textbf{Suricata} y \textbf{Snort 3}, complementadas con la referencia metodológica de \textbf{Wireshark}.
  \item \textbf{Dimensión Arquitectónica:} Modelado de una topología de red empresarial jerárquica y coherente, estableciendo los puntos de inserción para sensores perimetrales, sensores de distribución interna y sensores a nivel de host (\textit{Host-based Network IPS}, H-NIPS).
  \item \textbf{Dimensión de Implementación:} Despliegue de la solución nativa basada en \textbf{Suricata 7} en el servidor \texttt{dc1} (\texttt{192.168.0.2}), integrada con el cortafuegos \textbf{nftables} a través del módulo de colas \textbf{NFQUEUE}.
  \item \textbf{Dimensión de Reglas y Validación:} Construcción de firmas específicas para la protección de servicios de Active Directory (Kerberos en puerto 88, LDAP en puerto 389, SMB en puerto 445), fuerza bruta SSH (puerto 22) y ataques a servicios web (puertos 80/443), incluyendo la validación experimental mediante pruebas de detección y descarte activo en tiempo real.
\end{itemize}

\subsection{Límites y Restricciones}

\begin{itemize}
  \item \textbf{Entorno de Red:} El despliegue físico y empírico se encuentra acotado a la infraestructura del laboratorio en la subred plana \texttt{192.168.0.0/24}, bajo un servidor físico con Debian GNU/Linux y un conmutador/enrutador de borde TP-Link TL-WR850N.
  \item \textbf{Tráfico Cifrado de Extremo a Extremo:} Las capacidades de inspección de contenido profundo en protocolos con cifrado TLS 1.3 (HTTPS) se limitan al análisis de los metadatos de la negociación inicial (\textit{Server Name Indication}, SNI, y huellas digitales de certificados), dado que no se realiza ruptura forzada de certificados (\textit{TLS Interception / Man-in-the-Middle}) por motivos de privacidad y normativas de confidencialidad.
  \item \textbf{Rendimiento de Hardware:} Las pruebas de desempeño se ejecutan sobre hardware de laboratorio (servidor con procesador multinúcleo x86\_64 y memoria RAM de 8 GB), abstrayendo entornos corporativos de múltiples decenas de gigabits por segundo donde se requerirían aceleradores dedicados basados en FPGA o controladores con tecnología DPDK/SmartNIC.
\end{itemize}
